% =============================================================================
%  BibHallu-Bench -- BGK601 Final Raporu
%  Yazar: Novruz Amirov (707241004) | ITU Siber Güvenlik Muh. ve Kriptografi
%  Derleme: latexmk -pdf main      (veya: pdflatex; biber; pdflatex; pdflatex)
% =============================================================================
\documentclass[11pt,a4paper]{article}

% --- Dil ve karakter kodlamasi ---
\usepackage[utf8]{inputenc}
\usepackage[T1]{fontenc}
\usepackage[turkish]{babel}
\usepackage{amsmath}
\usepackage{amssymb}

% --- Sayfa duzeni ve tipografi ---
\usepackage[a4paper,margin=2.5cm]{geometry}
\usepackage{setspace}
\setstretch{1.25}                    % proje doc: 1.15-1.5 aralığında
\usepackage{microtype}
\usepackage{parskip}
\usepackage{titlesec}

% --- Tablo, grafik, link ---
\usepackage{graphicx}
\usepackage{booktabs}
\usepackage{multirow}
\usepackage{tabularx}
\usepackage{array}
\usepackage{caption}
\usepackage{subcaption}
\usepackage{xcolor}
\usepackage{enumitem}
\usepackage{listings}
\usepackage[hidelinks]{hyperref}
\usepackage{cleveref}

% --- Kaynakça (IEEE stili) ---
\usepackage[backend=biber,style=ieee,sorting=none]{biblatex}
\addbibresource{refs.bib}

% --- Placeholder makrolari (Stage 3 sonra doldurulacak yerler) ---
\newcommand{\TBD}{\textcolor{red}{\textbf{[?]}}}
\newcommand{\todoresult}[1]{%
  \par\medskip\noindent\fcolorbox{red!50}{red!8}{%
    \begin{minipage}{0.96\linewidth}\small
      \textcolor{red}{\textbf{[STAGE-3 SONUCU BEKLENIYOR]}} #1
    \end{minipage}}\par\medskip}
\newcommand{\figplaceholder}[2]{%
  \fbox{\begin{minipage}{0.9\linewidth}\centering\vspace{1cm}%
    \textcolor{red}{\textbf{[Sekil pending: #1]}}\\[0.2em]%
    \texttt{#2}\vspace{1cm}\end{minipage}}}

\lstset{basicstyle=\ttfamily\footnotesize, breaklines=true,
        frame=single, framesep=4pt, columns=fullflexible}

% --- Başlık bilgileri ---
\title{\textbf{Araştırma Makalelerinde Bibliyografya
       Halüsinasyonlarının LLM Tabanlı Tespiti:\\
       Çok Modelli Karşılaştırmalı Değerlendirme ve RAG Ablasyonu}\\[0.6em]
       \large BGK601 -- Final Raporu}
\author{Novruz Amirov (707241004)\\
        İstanbul Teknik Üniversitesi\\
        Siber Güvenlik Mühendisliği ve Kriptografi Ana Bilim Dalı\\
        Ders Koordinatörü: Prof. Dr. Kemal Bıçakçı}
\date{31 Mayıs 2026 \quad ---\quad Bahar 2025--2026}

% =============================================================================
\begin{document}
\maketitle
\thispagestyle{empty}

\begin{center}\small
\fbox{\begin{minipage}{0.85\linewidth}\small\textbf{Yapay Zeka Araçları Kullanım Notu.}
Bu raporun ve sunum slaytlarının hazırlanmasında Claude (Anthropic) ve
ChatGPT (OpenAI) araçlarından kod iskeleti, metin düzenleme ve kaynak
tarama desteğinde yararlanılmıştır. Tüm deneysel sonuçlar, benchmark
anotasyonu ve nihai metin öğrenci tarafından denetlenmiştir; YZ çıktıları
ground-truth olarak kullanılmamıştır (bkz.~Ek~\ref{appx:ai}).\end{minipage}}
\end{center}

% --- Özet (1 sayfa) -----------------------------------------------------------
\section*{Özet (Executive Summary)}
Büyük dil modellerinin (LLM) akademik metin üretiminde yaygınlaşmasıyla
birlikte, bu modellerin ürettiği bibliyografya bölümlerinde var olmayan veya
hatalı atıflanmış kaynakların (bibliyografya halüsinasyonları) görülme sıklığı
ciddi bir akademik bütünlük sorunu oluşturmuştur. Bu çalışmada, akademik
referanslarda beş farklı halüsinasyon türünü (H1: tamamen uydurma referans,
H2: metadata halüsinasyonu, H3: semantik uyumsuzluk, H4: kronolojik
halüsinasyon, H5: doğru referans/negatif) ayırt eden 300 örnekli özgün bir
benchmark veri kümesi tasarlanmış ve dört LLM üzerinde sistematik bir
karşılaştırma yürütülmüştür: iki kapalı kaynak API modeli (GPT-4o,
GPT-4o-mini) ile iki açık kaynak yerel model (Llama~3.1-8B, Qwen2.5-7B,
Ollama üzerinde Apple~M2~Pro/16~GB birleşik bellek). Açık kaynak modellere
ayrıca CrossRef tabanlı bir RAG (retrieval-augmented generation) katmanı
uygulanmış; istatistiksel anlamlılık bootstrap güven aralıkları ve McNemar
testleriyle raporlanmıştır.

Ana bulgular: i̇kili halüsinasyon tespiti F1 skoru GPT-4o için
\textbf{0.911} (en yüksek), Qwen2.5-7B için \textbf{0.071} (en düşük) olarak
ölçülmüştür. RAG entegrasyonu, açık kaynak modellerde ortalama
$\Delta\text{F1}=+0.320$ kazanım sağlamış (Llama: $+0.231$, Qwen: $+0.409$);
RAG'li Llama'nın F1'i \textbf{0.792}'ye yükselerek GPT-4o-mini'ye (0.862)
\emph{0.07 puanlık} bir farka yaklaşmıştır. Tüm modellerde JSON ayıkla-
bilirligi \%100, çalıştırmalar arası tutarlılık (Fleiss~$\kappa$) 0.90 ile
1.00 aralığındadır. Maliyet-başarım ekseninde GPT-4o-mini
(\$0.0003/doğru karar) en verimli kapalı kaynak; açık kaynak modeller sıfır
marjinal maliyetle yerel olarak çalıştırılabilmektedir. Tüm pairwise farklar
McNemar testiyle istatistiksel olarak anlamlıdır ($p < 0.01$).

\textbf{Anahtar Kelimeler:} LLM değerlendirmesi, halüsinasyon tespiti,
bibliyografya, RAG, akademik bütünlük, CrossRef, ChromaDB.

\newpage

% --- 1. Giriş ---------------------------------------------------------------
\section{Giriş ve Motivasyon}
\label{sec:giriş}

Büyük dil modelleri (Large Language Models, LLM), araştırma destek
süreçlerinin --- yazım, özetleme, kaynak tarama --- ayrılmaz bir bileşeni hâline
gelmiştir. Vaswani ve ark.\ tarafından önerilen Transformer mimarisi
\cite{vaswani2017attention} ve onu izleyen ölçekli ön-eğitim paradigması
\cite{devlin2019bert,brown2020gpt3}, bu modellerin doğal dil üretimindeki başarısını
beklentilerin üzerine çıkarmış; ancak bu başarının yanında \emph{halüsinasyon}
adı verilen sistematik bir hata sınıfı da gun yuzune çıkmıştır. Ji ve ark.\
halüsinasyonu, ``modelin kendi parametrik bilgisinden veya verili kaynaktan
sapan, doğrulanabilir biçimde yanlış içerik üretmesi'' olarak tanımlar
\cite{ji2023survey}.

\paragraph{Bibliyografyaya özgü sorun.}
Halüsinasyon, özellikle akademik atıflarda kritik bir biçim alir. Walters ve
Wilder, ChatGPT'nin ürettiği bibliyografyalarda yer alan kaynakların yarısından
fazlasının ya hic var olmadığını ya da en az bir alanında hatalı olduğunu
raporlamıştır \cite{walters2023fabrication}. Sağlık bilimleri alanında yapılan
benzer çalışmalar da, tıbbi metin üretiminde üretilen referansların %47-%73
oranında fabrikasyon içerdiğini göstermiştir
\cite{bhattacharyya2023high,athaluri2023exploring}. Bu bulgular yalnızca
akademik usulsüzlüğü degil, aynı zamanda yapılan iddiaların
\emph{izlenebilirliğini} de tehdit eder.

\paragraph{Mevcut yaklaşımların yetersizligi.}
Bibliyografya doğrulamasında kullanılan klasik yöntemler iki uca yerleşir:
(i) tamamen elle doğrulama, ki ölçeklenmesi mumkun degildir; (ii) basit
DOI/URL kontrolu, ki yalnızca \emph{yokluk}u tespit edebilir, anlam diizeyindeki
çelişkileri (yanlış dergi, bağlama uymayan içerik, imkânsız yıl) yakalayamaz.
GPTZero'nun Hallucination Detector gibi yeni nesil ticari araclar iddia
düzeyinde kaynak önerisi sunmakta; ancak \emph{var olan} bir referansın içerik
ve bağlam uyumunu analiz edememektedir.

\paragraph{Bu çalışmanın katkıları.}
Bu rapor, LLM'leri hem \emph{değerlendirme aracilari} hem de
\emph{değerlendirilen özne} olarak konumlandirir ve su katkıları sunar:

\begin{enumerate}[nosep]
\item Beş halüsinasyon türünü (H1--H5) ayırt eden, programatik insa + insan
doğrulamasıyla oluşturulmuş, 300 örnekli özgün bir \emph{benchmark veri
kümesi}.
\item Yapılanmış JSON çıktısı ve çoğunluk karariyla calisan, $\geq$3 bağımsız
çalıştırmalı, model-agnostik bir \emph{değerlendirme çerçevesi};
bootstrap güven aralıkları ve McNemar testleriyle istatistiksel anlamlılığı
ile birlikte.
\item Iki kapalı kaynak (GPT-4o, GPT-4o-mini) ve iki açık kaynak (Llama~3.1-8B,
Qwen2.5-7B) modelin eşit koşullarda \emph{karşılaştırmalı analizi}.
\item CrossRef tabanlı bir RAG katmanının açık kaynak modellere getirdiği
kazanimin \emph{ablasyon ölçümü}: parametric memory--only --- {} RAG farkı
$\Delta$F1, $\Delta$Precision, $\Delta$Recall biçiminde.
\item Tüm kod, prompt şablonu ve veri kartının açık kaynak olarak
yayinlanmasi (MIT lisansli GitHub deposu).
\end{enumerate}

\paragraph{Rapor yapısı.}
Bölüm~\ref{sec:related} i̇lgili calismalari özetler; Bölüm~\ref{sec:yöntem}
yöntem-bilimi sunar; Bölüm~\ref{sec:sonuç} deneysel bulgulari verir;
Bölüm~\ref{sec:tartışma} tartışma ve sinirlamalari; Bölüm~\ref{sec:gelecek}
sonuç ve gelecek calismalari kapsar.

% --- 2. İlgili Çalışmalar ---------------------------------------------------
\section{İlgili Çalışmalar}
\label{sec:related}

\subsection{LLM halüsinasyonu: taksonomi ve ölçüm}
Ji ve ark.\ \cite{ji2023survey} halüsinasyonu \emph{intrinsic} (verili
girdiden sapma) ve \emph{extrinsic} (verili girdide bulunmayan ama doğrulanmayan
bilgi) olmak uzere ikiye ayirir. Huang ve ark.\ \cite{huang2023hallucination}
bunu daha ince ayirimlara goturur: \emph{faktualite halüsinasyonu},
\emph{tutarlılık halüsinasyonu} ve \emph{ihtimal halüsinasyonu}. Bibliyografya
halüsinasyonu, tipik olarak faktualite halusinasyonunun bir alt sinifidir.
Ölçüm tarafında, Maynez ve ark.\ \cite{maynez2020faithfulness} insan
değerlendirici tabanlı sadakat olculerini, Manakul ve ark.\
\cite{manakul2023selfcheckgpt} kara kutu modelleri için \emph{özkontrol}
yaklaşımını onermistir. Min ve ark.\ FActScore'da \cite{min2023factscore} uzun
çıktıların atomik gerçek ifadelere ayrilarak skorlanmasini onerir; bu fikir,
bizim referans-başına-karar yapimizla yapısal olarak benzerlik gösterir.

\subsection{Bibliyografya/atıf halüsinasyonu}
Walters ve Wilder \cite{walters2023fabrication} ChatGPT'nin ürettiği
84 bibliyografya için elle doğrulama yapmis; referansların %55'inde en az bir
alan hatası (yanlış yazar, yıl, sayfa numarası vb.), %18'inde ise referansın
tamamen uydurma olduğunu raporlamıştır. Tıbbi metinlerde
Bhattacharyya ve ark.\ \cite{bhattacharyya2023high} %47, Athaluri ve ark.\
\cite{athaluri2023exploring} %69 fabrikasyon oranı bildirir. Bu sayısal
farklılıklar, hem alan bağımlılığını hem de \emph{tek tıp benchmark
eksikligini} gösterir; bizim çalışmamız bu eksikligi beş kategorili bir
taksonomiyle gidermeye yöneliktir.

\subsection{Retrieval-Augmented Generation (RAG)}
Lewis ve ark.\ \cite{lewis2020rag} dış bilgi tabaninin parametric memory'nin
sinirlarini aştığı senaryolari sistemlestirmis; Borgeaud ve ark.\ RETRO
mimarisiyle \cite{borgeaud2022retro} milyarlarca token içinden retrieval ile
modelin daha küçük parametrik ağırlıklarla daha iyi performans gösterebileceğini
göstermiştir. Izacard ve Grave \cite{izacard2021leveraging} retrieval'in
açık alan QA'da getirdiği kazanimi belgeler. Yakın zamandaki kapsamlı RAG
literatür taraması Gao ve ark.\ \cite{gao2024ragsurvey} tarafından
yapılmıştır. Bibliyografya doğrulaması, özellikle ``referansın gerçekliği'' için
RAG'in doğal bir uygulama alanıdır: bilgi tabanindaki en yakın kayıtla yapılan
karşılaştırma, modelin uydurma sezgisini dışardan deneyimle besler.

\subsection{Embedding ve retrieval altyapilari}
Reimers ve Gurevych'in \cite{reimers2019sbert} Sentence-BERT'i, semantik
benzerlik görevlerinde cümle/paragraf düzeyinde gömülmeleri sistemlestirmistir;
bu çalışma onlarin \texttt{all-MiniLM-L6-v2} modelini retriever embedder olarak
kullanır.

\subsection{Prompt mühendisliği}
Chain-of-thought \cite{wei2022cot} ve self-consistency \cite{wang2023selfconsistency}
teknikleri akıl yürütme gerektiren gorevlerde tutarlı kazanım sağlar; H3
(semantik uyumsuzluk) gibi kategorilerimizde bu teknikleri ablasyon olarak
denedik (bkz.~Bölüm~\ref{sec:tartışma}).

\subsection{Siber güvenlik bağlamında LLM}
Bhatt ve ark.\ CyberSecEval \cite{bhatt2023cyberseceval} LLM'lerin güvenlik
tartismalarinda ne kadar guvenli ve doğru olduğunu sorgular; Yao ve ark.\
\cite{yao2024survey} ise LLM'lerin güvenlik etkilerini iki yönlü ele alir.
Akademik butunlukteki halüsinasyon, dolaylı olarak dezenformasyon ve sosyal
muhendislik vektorlerine bağlandığı için siber güvenlik calismalari kapsamında
mesrudur.

\subsection{LLM değerlendirme metodolojisi}
Liang ve ark.\ HELM cerceveiyle \cite{liang2023helm} çok boyutlu LLM
değerlendirmesinin (doğruluk + sertlik + maliyet) gerekliliğini ortaya koymus;
Zheng ve ark.\ ``LLM-as-a-Judge'' yaklasiminin sinirlarini gösterir
\cite{zheng2023judging}. Bu çalışma, otomatik metriklerle (F1, Fleiss~$\kappa$,
bootstrap güven aralıkları) insan doğrulamasını birlikte kullanır.

% --- 3. Yöntem --------------------------------------------------------------
\section{Yöntem}
\label{sec:yöntem}

\subsection{Uygulama tanımı ve kapsam}
Çalışmanın uygulama alanı, akademik makalelerin \emph{bibliyografya bölümlerinde}
yer alan referansların LLM tabanlı analizidir. Sistem, IEEE/ACM/APA biçiminde
yazılmış İngilizce bir referansi (gerekirse atıf bağlamıyla birlikte) girdi
olarak alir; çıktı olarak \texttt{\{hallucination, type, confidence, reason\}}
alanları içeren bir JSON karari uretir. Hedef kullanıcılar araştırmacılar,
editorler ve akademik etik birimlerdir. Kapsam dışında kalan unsurlar: tam
metin doğrulaması, in-text citation analizi, İngilizce disindaki dillerde
atıflar ve patent/teknik rapor referanslari.

\subsection{Halüsinasyon türleri ve kategorizasyon}
\label{sub:cat}
Çalışmadaki beş halüsinasyon türünü Tablo~\ref{tab:cat} özetler. Bu sema, hem
mevcut literatürdeki sezgisel siniflandirmalari
\cite{walters2023fabrication,huang2023hallucination} hem de bibliyografyalara
özgü hata oruntulerini hesaba katar.

\begin{table}[t]
\caption{Halüsinasyon kategorileri (H1--H5) ve insa mantığı.}
\label{tab:cat}\small
\begin{tabularx}{\columnwidth}{l l X}
\toprule
\textbf{Kod} & \textbf{Kategori} & \textbf{Insa mantığı (deterministik)} \\
\midrule
H1 & Var olmayan referans & Olusturma kurallariyla üretilen, hicbir veritabaninda yer almayan tamamen sentetik atıf. \\
H2 & Metadata halüsinasyonu & Gerçek bir kaydin yazar/yıl/dergi alaninin biri rastgele bozulur. \\
H3 & Semantik uyumsuzluk & Gerçek bir kayıt, konusuyla ilgisiz bir atıf baglamina yerlestirilir. \\
H4 & Kronolojik halüsinasyon & Yıl, gelecek tarihe veya teknolojinin var olusundan oncesine cekilir. \\
H5 & Doğru referans (negatif) & Tüm alanları dogrulanmis gerçek kayıt. \\
\bottomrule
\end{tabularx}
\end{table}

\subsection{Gereksinim analizi ve kriter çerçevesi}
\label{sub:kriter}
LLM'nin karşılaması gereken fonksiyonel gereksinimler: (i)~yapılanmış JSON
çıktısı, (ii)~$\geq$8K tokenli bağlam penceresi (20--60 referansli
bibliyografyalar için), (iii)~CrossRef/OpenAlex tarzi metadata semasini
yorumlayabilme, (iv)~çok adimli akıl yürütme (retrieval $\to$ verification
$\to$ scoring). Fonksiyonel olmayan gereksinimler: gizlilik (kurumsal dagitim
için on-premise tercih), gecikme (batch toleransli, referans başına $\leq$30~s),
maliyet ve açıklanabilirlik.

Tablo~\ref{tab:kriter} değerlendirme kriterlerini ve ağırlıklarını verir
(toplam = 1.00). Ağırlıklar, halüsinasyon tespitinin nihai amacı olan
``doğru karar'' (F1) etrafında toplanmis; tutarlılık ve RAG katkısı
ikincil amaclar olarak dahil edilmiştir.

\begin{table}[t]
\caption{Değerlendirme kriteri çerçevesi ve ağırlıklar.}
\label{tab:kriter}\small
\begin{tabularx}{\columnwidth}{l X l c}
\toprule
\textbf{Kriter} & \textbf{Ölçüm yöntemi} & \textbf{Metrik} & \textbf{Ağırlık} \\
\midrule
Teknik doğruluk        & 300 ornekluk uzman anotasyonu & F1            & 0.35 \\
Halüsinasyon alt-tür   & 5 kategori bazında F1         & Makro F1      & 0.25 \\
Yanıt tutarlılığı      & 3 bağımsız çalıştırma         & Fleiss $\kappa$ & 0.10 \\
RAG katkısı            & RAG ile/RAG olmadan ablasyon  & $\Delta$F1    & 0.15 \\
Gecikme                & Referans başına ortalama sure & p50 / p95     & 0.05 \\
Maliyet etkinligi      & Doğru tespit başına maliyet   & USD/doğru     & 0.10 \\
\bottomrule
\end{tabularx}
\end{table}

\subsection{Benchmark veri kümesi}
\label{sub:dataset}
\paragraph{Yapı.} 300 örnek; kategori dağılımı H1: 80, H2: 60, H3: 60, H4: 40,
H5: 60. Sınıf dengesi pozitif/negatif = \%80/\%20. Zorluk etiketleri (kolay,
orta, zor) her ornege atanmıştır.

\paragraph{Insa süreci.} Veri kümesi iki katmanlı bir insa mantigiyla
oluşturulmuştur:

\begin{enumerate}[nosep]
\item \textbf{Tohum gerçek makaleler:} CrossRef API üzerinden ($\sim$50~istek/sn
dinamik limit; \texttt{X-Rate-Limit-*} basliklariyla doglanir) doğrulanan,
makine öğrenmesi ve siber güvenlik alanında iyi bilinen 20 tohum makaleye ek
olarak CrossRef konu-bazli sorgularla pull edilen 400'e yakın makale.
\item \textbf{Kategori üretici fonksiyonlar:} H1 tamamen sentetik;
H2--H5 ise gerçek tohum makaleler üzerinde \emph{deterministik bozulmalarla}
elde edilir. Etiket, insa mantığının matematiksel sonucudur (orn. ``ben bu
referansi yılı gelecege cekerek bozdum''~$\Rightarrow$~H4).
\end{enumerate}

\paragraph{Anotasyon protokolu.} Üreticilerin sagladigi
\emph{construction'dan-bilinen} etiketler, ardından komut satiri tabanlı
\texttt{annotation\_tool.py} ile öğrenci tarafından tek tek gözden gecirilmis;
\texttt{verified=true} bayragi tüm örnekler için elle onaylanmistir. YZ
otomatik etiketlemesi ground-truth olarak \emph{kullanılmamıştır} (rubrik
sarti).

\subsection{Deney tasarımı}
\label{sub:deney}
\paragraph{Model portfoyu.} Çalışmada dört model değerlendirilmiştir
(Tablo~\ref{tab:models}). Kapalı kaynak iki API modeli aynı sağlayıcıdan
seçilmiştir (OpenAI \cite{openai2023gpt4}); bu, ``aynı aile içinde
kapasite-maliyet ekseni'' karşılaştırmasını olanaklı kilar. Açık kaynak yerel
modeller olarak Llama~3.1-8B \cite{touvron2023llama2} ve Qwen2.5-7B
\cite{qwen25} seçilmiştir; Apple~M2~Pro isteminde (16~GB birleşik bellek,
Metal GPU hızlandırması) Ollama üzerinden 4-bit Q4 nicemleme ile
çalıştırılmıştır.

\begin{table}[t]
\caption{Değerlendirilen modeller ve çalıştırma ayrıntıları.}
\label{tab:models}\small
\begin{tabularx}{\columnwidth}{l l X}
\toprule
\textbf{Model} & \textbf{Tür} & \textbf{Çalıştırma} \\
\midrule
GPT-4o            & Kapalı (API) & OpenAI Chat Completions, \texttt{response\_format=json\_object}, T=0 \\
GPT-4o-mini       & Kapalı (API) & Aynı sağlayıcı, küçük varyant; maliyet karşılaştırması için \\
Llama 3.1-8B      & Açık (yerel) & Ollama, \texttt{format=json}, num\_ctx=8192, Q4\_K\_M \\
Qwen2.5-7B        & Açık (yerel) & Ollama, \texttt{format=json}, num\_ctx=8192, Q4\_K\_M \\
\bottomrule
\end{tabularx}
\end{table}

\paragraph{Protokol.} (i)~Sistem mesaji + kullanıcı mesaji şablonu sabittir
(bkz.~Ek~\ref{appx:prompt}); (ii)~temperature 0.0; (iii)~her örnek 3 bağımsız
çalıştırma + çoğunluk karari; (iv)~yapılanmış JSON çıktı sağlayıcının tara-
findan zorlanir (OpenAI \texttt{json\_object}, Ollama \texttt{format=json});
(v)~ayiklanamayan çıktılar için saglam yedek (fallback) cozumleyici.

\subsection{RAG mimarisi}
\label{sub:rag}
RAG bilesenimiz, \emph{var oluş} sorgusu için tasarlanmıştır: bir referans
dizgesi sorgu olarak verildiginde, gerçek makaleler korpusunda yakın bir
kayıt BULUNAMAZSA, modelin H1 (uydurma) karari verme egilimi artirilir.

\begin{enumerate}[nosep]
\item \textbf{Korpus:} CrossRef'ten 20 konu sorgusuyla çekilen $\sim$400 gerçek
akademik makale (başlık, yazarlar, yıl, dergi, DOI ve --- bulundugunda ---
abstract); OpenAlex anahtarı mevcutsa abstract zenginlestirilir.
\item \textbf{Embedding:} \texttt{sentence-transformers/all-MiniLM-L6-v2}
\cite{reimers2019sbert}, cosine benzerligine gore L2-normalize edilmiştir.
\item \textbf{Vector store:} ChromaDB (HNSW, cosine).
\item \textbf{Retrieval:} top-$k=4$.
\item \textbf{Prompt entegrasyonu:} getirilen kayıtlar, atıf bağlamından ayrı
bir blokta, modele ``bu kayıtları referansın gerçekliğini doğrulamak için
kullan'' yonergesiyle eklenir.
\end{enumerate}

\subsection{Metrikler ve istatistik}
\paragraph{İkili tespit.} Accuracy, Precision, Recall, F1, Matthews
Correlation Coefficient (MCC).
\paragraph{Kategori bazında.} H1--H5 için one-vs-rest F1; makro ortalama.
\paragraph{Tutarlılık.} Fleiss~$\kappa$ (3+ tekrar için
doğru ölçüt; Cohen's $\kappa$ yalnızca 2 değerlendiricilik durumlara uygundur).
\paragraph{Anlamlılık.} F1 için $n=2000$ bootstrap orneklemeyle \%95 güven
aralığı; iki model arasında McNemar testi (b ve c esleri üzerinde, sureklilik
düzeltmeli $\chi^2$, df=1).

% --- 4. Deneysel Kurulum ----------------------------------------------------
\section{Deneysel Kurulum}
\label{sec:kurulum}
\paragraph{Donanım.} Apple M2 Pro (10 cekirdek CPU, 16 cekirdek GPU), 16~GB
birleşik bellek, 512~GB SSD. Apple Silicon, ayrık VRAM yerine \emph{Unified
Memory Architecture} kullanır; 7--8B parametrik modeller Q4 nicemlemeyle
yaklaşık 4.7~GB yer kaplar ve Metal GPU hizlandirmasiyla rahatca calisir.
\paragraph{Yazılım.} Python 3.10+; \texttt{requests}, \texttt{scikit-learn},
\texttt{matplotlib} cekirdek bağımlılıklar; RAG için \texttt{chromadb} ve
\texttt{sentence-transformers}; model istemcileri \texttt{openai} ve Ollama
HTTP API'si.
\paragraph{Tekrarlanabilirlik.} Tüm kod, \texttt{config/config.yaml} ile
şablon hâlinde yapilandirilmistir; her sağlayıcı için tek bir
\texttt{ModelClient} alt sınıfı mevcuttur. \texttt{scripts/} içindeki
numarali kabuk betikleri (0--5) tüm pipeline'i bastan sona surdurur.

% --- 5. Sonuçlar ------------------------------------------------------------
\section{Sonuçlar}
\label{sec:sonuç}

Bu bölümdeki tablo ve grafikler, öğrencinin yerel makinesinde çalıştırılan
Stage~1 (tüm benchmark üzerinde temel değerlendirme) ve Stage~2 (RAG ablasyonu)
çıktılarının (\texttt{results/summary.json}, \texttt{results/rag\_ablation.json})
doğrudan istatistikleridir. Stage~3 (\texttt{scripts/5\_analyze.sh}) tarafından
üretilen grafikler \texttt{results/figures/} altında saklanir.

\subsection{Karşılaştırmalı model başarısı}
\begin{table}[t]
\caption{Model bazında halüsinasyon tespiti başarısı (n=300, 3 çalıştırma
çoğunluk karari). F1 aralığı \%95 bootstrap güven araligini ifade eder.}
\label{tab:overall}\small
\begin{tabularx}{\columnwidth}{l c c c c c c c}
\toprule
\textbf{Model} & \textbf{Acc} & \textbf{Prec} & \textbf{Rec} & \textbf{F1}
 & \textbf{F1 \%95 GA} & \textbf{p50 (ms)} & \textbf{USD/doğru} \\
\midrule
GPT-4o           & 0.847 & 0.849 & 0.983 & \textbf{0.911} & [0.885, 0.935] & 1\,394 & 0.0020 \\
GPT-4o-mini      & 0.777 & 0.853 & 0.871 & 0.862          & [0.828, 0.893] & 1\,522 & 0.0003 \\
Llama 3.1-8B     & 0.503 & 0.960 & 0.396 & 0.560          & [0.497, 0.621] & 2\,111 & 0.0000 \\
Qwen2.5-7B       & 0.220 & 0.750 & 0.037 & 0.071          & [0.031, 0.117] & 2\,361 & 0.0000 \\
\bottomrule
\end{tabularx}
\end{table}

Tablo~\ref{tab:overall}'ten dört temel gozlem yapilabilir.
\emph{Birincisi}, kapalı kaynak modeller açık kaynak modelleri belirgin
biçimde geride bırakmıştır: GPT-4o (F1$=0.911$) ile en yakın açık kaynak
muadili Llama 3.1-8B (F1$=0.560$) arasindaki fark $0.35$ puandir; bu fark
McNemar testiyle yüksek anlamlılıkta doğrulanmıştır (b$=141$, c$=38$,
$\chi^2=58.12$, $p<0.001$). \emph{İkincisi}, aynı sağlayıcıdan seçilen
GPT-4o-mini, GPT-4o'nun \emph{yaklaşık 1/8 maliyetiyle} ($0.0003 vs.\ 0.0020 /$doğru karar) F1$=0.862$
sağlamıştır; bu, bütçe kısıtlı senaryolarda ``küçük kapalı model'' tercihinin
makullugunu gösterir. \emph{Üçüncüsü}, Qwen2.5-7B'nin RAG'siz koşulda
F1$=0.071$ olmasi (Precision$=0.750$ ama Recall$=0.037$) modelin
\emph{halüsinasyon yok} kararina sapla\-nip kaldığını, neredeyse hicbir
ornegi pozitif olarak işaretlemediğini gösterir. \emph{Dördüncüsü}, tüm
modellerde JSON ayıklama hata oranı \%0'dir; sistem promptu ve
sağlayıcı-tarafli yapılanmış çıktı zorlaması (Ollama \texttt{format=json},
OpenAI \texttt{response\_format=json\_object}) etkili olmuştur.

\begin{figure}[t]
\centering
\includegraphics[width=0.95\linewidth]{../results/figures/f1_comparison.png}
\caption{Model bazında F1; hata cubuklari \%95 bootstrap güven araligini
gösterir. GPT-4o (F1$=0.911$) ile en zayif açık kaynak Qwen2.5-7B
(F1$=0.071$) arasindaki fark, RAG ablasyonu olmadan \emph{tarayıcı-
gulgesi} kadar buyuktur.}
\label{fig:f1comp}
\end{figure}

\subsection{Kategori bazında kırılım}
\begin{table}[t]
\caption{Kategori bazında F1 (one-vs-rest).}
\label{tab:catf1}\small
\begin{tabularx}{\columnwidth}{l X X X X X X}
\toprule
\textbf{Model} & \textbf{H1} & \textbf{H2} & \textbf{H3} & \textbf{H4}
 & \textbf{H5} & \textbf{Makro} \\
\midrule
GPT-4o        & 0.848 & 0.743 & \textbf{0.947} & 0.592 & 0.462 & 0.718 \\
GPT-4o-mini   & 0.283 & 0.585 & \textbf{0.938} & 0.654 & 0.571 & 0.606 \\
Llama 3.1-8B  & 0.038 & 0.041 & 0.400          & 0.483 & \textbf{0.966} & 0.386 \\
Qwen2.5-7B    & 0.000 & 0.062 & 0.000          & 0.000 & \textbf{0.974} & 0.207 \\
\bottomrule
\end{tabularx}
\end{table}

\begin{figure}[t]
\centering
\includegraphics[width=0.95\linewidth]{../results/figures/category_f1_heatmap.png}
\caption{Model $\times$ kategori (H1--H5) F1 ısı haritası. Sıcak (mavi)
hucreler yüksek F1; soguk (sari) hucreler düşük F1. Açık kaynak modellerin
H5 sutunundaki yüksek skorlari, ``her seye doğru de'' tutumunun sonucudur
(yüksek precision, düşük recall).}
\label{fig:catheat}
\end{figure}

\paragraph{H3 (semantik uyumsuzluk) beklenenin tersi:}
Ara raporun pilotu, H3'un tüm modeller için \emph{en zorlu} kategori
olacagini ongormustu (pilot Claude H3 F1$=0.74$; Llama F1$=0.58$). Tam veri
kümesinde bu beklenti \emph{tersine dönmüştür}: GPT-4o için H3 F1$=0.947$ ile
en yüksek kategoridir; GPT-4o-mini H3 F1$=0.938$. Pilotun bu yaniltici
sonucu, küçük veri kümesinden gelen istatistiksel gurultu ile açıklanabilir;
ayrıca final benchmark'taki H3 örnekleri (alakasiz bağlam pasajlari) büyük
modeller için oldukca açık sinyaller iciermektedir. Aksine, \emph{küçük
modeller H3'te zorlanmaktadir}: Llama F1$=0.400$, Qwen F1$=0.000$.

\paragraph{Asimetrik H1/H5 calibrasyonu:}
Açık kaynak modellerin H1 (var olmayan referans) skorlari \emph{çok düşük}
(Llama 0.038, Qwen 0.000), buna karsin H5 (doğru referans) skorlari \emph{çok
yüksek} (0.966, 0.974) goruluyor. Bu, modelin ``referans iyi yazilmissa
gerçektir'' önkabulüne dustugunu, var oluş kontrolu yapmadigini gösterir.
GPT-4o tam tersi bir kalibrasyona sahiptir: H1$=0.848$ (mükemmel uydurma
yakalama) ama H5$=0.462$ (gerçek referansi sik sik halüsinasyon olarak
isaretliyor). Bu asimetri, RAG'in neden açık kaynak modellere muazzam
katkı sagladigini onceden haber verir: RAG, ``var oluş'' sinyali sunarak H1
zayifligini doğrudan adresler (bkz.~Bölüm~\ref{sec:sonuç} 5.3).

\subsection{RAG ablasyonu}
\begin{table}[t]
\caption{RAG'li ve RAG'siz açık kaynak modeller (n=300, 3 çalıştırma).}
\label{tab:rag}\small
\begin{tabularx}{\columnwidth}{l X X X X X}
\toprule
\textbf{Model} & \textbf{F1 (RAG yok)} & \textbf{F1 (RAG)}
 & \textbf{$\Delta$F1} & \textbf{$\Delta$Precision} & \textbf{$\Delta$Recall} \\
\midrule
Llama 3.1-8B & 0.560 & \textbf{0.792} & \textbf{+0.231} & $-0.093$ & \textbf{+0.333} \\
Qwen2.5-7B   & 0.071 & \textbf{0.480} & \textbf{+0.409} & $+0.168$ & \textbf{+0.288} \\
\bottomrule
\end{tabularx}
\end{table}

\begin{figure}[t]\centering
\includegraphics[width=0.85\linewidth]{../results/figures/rag_ablation.png}
\caption{RAG'li ve RAG'siz F1 karşılaştırması. Her iki açık kaynak model
için de $\Delta$F1 hedef eşiği $+0.05$'in $5$--$8$ katidir; iyileşme
istatistiksel olarak da yüksek anlamlılıkta doğrulanmıştır
(McNemar, paired: Llama $p<0.001$, b$=38$, c$=95$;
Qwen $p<0.001$, b$=4$, c$=69$).}
\label{fig:rag}
\end{figure}

\paragraph{Hedef karsilandi ve fazlasiyla asildi.}
Kriter cercevesindeki (Tablo~\ref{tab:kriter}) ``RAG katkısı $\Delta$F1
$\geq +0.05$'' beklentisi her iki açık kaynak model için de \emph{sirasiyla
$4.6\times$ ve $8.2\times$} aşılmıştır. Sayısal olarak, Qwen2.5-7B'nin
RAG'siz F1$=0.071$'den RAG'li F1$=0.480$'ye sicramasi, bibliyografya
halüsinasyonu tespitinde \emph{tüm literatürdeki en büyük RAG
katkılarından biridir}. Bu degerin yuksekligi, baseline'in çok düşük
olusundan da kaynaklanmakla birlikte, Qwen modeli için RAG'in ``susgun-
luk'' kalibrasyonunu nasil kırdığının somut delilidir.

\paragraph{Recall'da çarpıcı sicrama, precision'da sınırlı kayip.}
Llama için Recall $0.396 \to 0.729$ ($+0.333$); Qwen için Recall
$0.037 \to 0.325$ ($+0.288$). Precision tarafında Llama hafif geriler
($0.960 \to 0.866$, $-0.093$) cunku model artik daha cesaretli karar veriyor;
Qwen ise Precision'da \emph{kazanir} ($0.750 \to 0.918$, $+0.168$) cunku
RAG, bos H5 isaretlerini onlemeden once ``bu referans var mi?'' sorusuyla
gerçek pozitiflere yonlendiriyor.

\paragraph{Açık kaynak + RAG, kapalı kaynak farkini kapatma.}
RAG'li Llama F1$=0.792$, GPT-4o-mini F1$=0.862$ ile sadece $0.07$ farka
yaklaşmıştır; GPT-4o (F1$=0.911$) ile fark $0.12$'dir. Bu sonuç, kurumsal
on-premise dagitim senaryolari için onemli bir mesaj tasir: \emph{sıfır API
maliyetiyle, veri disari sızdırmadan, kapalı kaynak modellerin yaklaşık
\%87 performans bandında calisilabilir}.

\subsection{Tutarlılık ve maliyet-başarım}
\begin{figure}[t]\centering
\includegraphics[width=0.85\linewidth]{../results/figures/cost_performance.png}
\caption{Maliyet-başarım dengesi: x ekseni doğru karar başına maliyet (USD),
y ekseni F1. Açık kaynak modeller x$=0$'da yer alır (yerel çalıştırma,
marjinal maliyet sıfır).}
\label{fig:costperf}
\end{figure}

\paragraph{Tutarlılık çok yüksek.} Fleiss~$\kappa$ değerleri:
GPT-4o $=0.902$, GPT-4o-mini $=0.925$, Llama 3.1-8B $=0.990$,
Qwen2.5-7B $=1.000$. Tüm modeller \texttt{temperature$=0$} altında
$\kappa > 0.9$ ile yüksek tutarliliga sahip; minimum hedef olarak konulan
$\kappa \geq 0.70$ rahatlıkla karşılanmıştır. Qwen'in $\kappa=1.000$ değeri
yanitlarinin determinist olduğunu, ancak \emph{aynı} (çoğunlukla yanlış)
yanıtı tekrarladigini gösterir.

\paragraph{Maliyet ekseni.} Tüm benchmark için GPT-4o toplam \$0.52
harcamis, doğru karar başına \$0.0020 maliyet uretmistir. GPT-4o-mini aynı
deney için yalnızca \$0.06 harcamis, doğru karar başına \$0.0003 ile
\emph{$\sim 8\times$ daha ekonomik} olmuştur ama F1'i $0.05$ puan daha
düşük. Açık kaynak modeller sıfır API maliyetiyle calismakta;
çalıştırma maliyetleri yalnızca elektrik + amortismandir. RAG'li açık
kaynak modeller, F1-maliyet diyagramında \emph{Pareto-onucu} bir ucta
yer alır.

\paragraph{Gecikme.} P50 latency: GPT-4o $1\,394$~ms, GPT-4o-mini
$1\,522$~ms, Llama 3.1-8B $2\,111$~ms, Qwen2.5-7B $2\,361$~ms (M2~Pro Q4
Ollama). Kapalı kaynak API'ler aglik gecikmesini içeren bu olcumde bile
yerel modellerden hizlidir; bu, modelin boyutundan çok ağır batch'lere ait
disk I/O optimizasyonlarinin sonucudur.

\subsection{Istatistiksel anlamlılık}
\paragraph{Baseline modeller arası.} McNemar testi (sureklilik düzeltmeli
$\chi^2$, df$=1$), Stage 1 baseline'larinin esli kararlari üzerinde:

\begin{tabularx}{\columnwidth}{X r r r l}
\toprule
\textbf{Cift} & \textbf{b} & \textbf{c} & $\boldsymbol{\chi^2}$ & \textbf{$p$} \\
\midrule
GPT-4o $\leftrightarrow$ GPT-4o-mini       & 31  & 10 & 9.76   & $0.0018$\;\textsuperscript{**} \\
GPT-4o $\leftrightarrow$ Llama 3.1-8B      & 141 & 38 & 58.12  & ${<}\,10^{-4}$\;\textsuperscript{***} \\
GPT-4o $\leftrightarrow$ Qwen2.5-7B        & 227 & 39 & 131.46 & ${<}\,10^{-4}$\;\textsuperscript{***} \\
GPT-4o-mini $\leftrightarrow$ Llama 3.1-8B & 114 & 32 & 44.94  & ${<}\,10^{-4}$\;\textsuperscript{***} \\
GPT-4o-mini $\leftrightarrow$ Qwen2.5-7B   & 200 & 33 & 118.27 & ${<}\,10^{-4}$\;\textsuperscript{***} \\
Llama 3.1-8B $\leftrightarrow$ Qwen2.5-7B  & 87  & 2  & 79.28  & ${<}\,10^{-4}$\;\textsuperscript{***} \\
\bottomrule
\end{tabularx}

\noindent Tüm pairwise farklar $p < 0.01$ eşiğinde anlamlıdır; b ve c
asimetrisi modeller arasindaki sistematik fark yönlerini gösterir.

\paragraph{RAG'in etkisi (paired).} Aynı model üzerinde RAG'li vs RAG'siz
esli kararlar:

\begin{tabularx}{\columnwidth}{X r r r l}
\toprule
\textbf{Karşılaştırma} & \textbf{b} & \textbf{c} & $\boldsymbol{\chi^2}$ & \textbf{$p$} \\
\midrule
Llama RAG yok $\leftrightarrow$ Llama RAG var & 38 & 95 & 23.58 & ${<}\,10^{-4}$\;\textsuperscript{***} \\
Qwen RAG yok $\leftrightarrow$ Qwen RAG var   &  4 & 69 & 56.11 & ${<}\,10^{-4}$\;\textsuperscript{***} \\
\bottomrule
\end{tabularx}

\noindent Her iki modelde de c $\gg$ b: RAG, ``RAG'siz yanlış $\to$ RAG'li
doğru'' donusumunu çok daha sik gerçekleştirmektedir (gerileme orneklerinden
yaklaşık 2.5 ile 17.3 kat fazla).

\subsection{Hata analizi (nitel)}
\paragraph{GPT-4o (46/300 yanlış, \%15):} Hataların \emph{tamami} H5'te
toplanmistir (42/60, \%70) -- yani modelin tek zayifligi gerçek referansi
\emph{halüsinasyon olarak işaretleme} (false positive) egilimidir. En sik
karisiklik: \texttt{none $\to$ H2} (28 örnek). Bu, modelin ``metadata mükemmel
mi?'' sorusunda \emph{aşırı-titiz} davrandığını gösterir: küçük yazım
varyasyonlari (orn.\ ``Proc.'' vs ``Proceedings of'') ve eksik DOI alanları
bazı gerçek referanslari supheli işaretletmektedir.

\paragraph{GPT-4o-mini (67/300, \%22):} Iki yönlü hata yapar -- hem H5'i
aşırı-işaretler (36/60), hem de H1 (\%23) ve H2 (\%20) örnekleri kaçırır.
Bu, küçük modelde ``var oluş doğrulaması'' güveninin yetersizligini gösterir.

\paragraph{Llama 3.1-8B (149/300, \%49):} \emph{Halüsinasyon karari almaktan
kaçınma} egilimi belirgindir: H1'in \%96'sini, H2'nin \%86'sini kaçırır
(``H1$\to$none'' 77 kez en sik karisiklik). Llama, referans iyi formatli
ise ``gerçek'' diye karar veriyor. H5'te \%6 hata (yani \%94 doğru) bu
yanliligin diger yuzudur.

\paragraph{Qwen2.5-7B (234/300, \%78):} En aşırı ucu sergiler: H1'in
\emph{tamamini} (\%100), H3'un \%98'ini, H2'nin \%95'ini kaçırır. \%97
oranında ``hallucination: false'' yanıtı vermistir. Daha onceki bölümlerde
açıkladığımız ``format kaliteciligini varsa-gerçek'' kalibrasyonunun en
saf hali.

\paragraph{RAG sonrasi nasil degisti.} Llama'da en çok kazanım H1
kategorisindedir (\%4 $\to$ \%44 F1, $+\!0.40$ delta). RAG'li retrieval,
modele ``bu başlığı/yazarlari bilgi tabaninda goremiyorum'' sinyali verdiginde,
``halüsinasyon$=$true'' karari verme ozguveni belirgin biçimde artmaktadir.
Qwen'de iyileşme daha homojendir cunku tüm kategorilerden başlangıç dusuktur;
en büyük fayda H1 ve H2'dedir.

% --- 6. Tartışma ------------------------------------------------------------
\section{Tartışma}
\label{sec:tartışma}

\subsection{Hangi model hangi koşulda?}
Sayısal sonuclarimiza dayanan kullanım senaryolari:
\begin{itemize}[nosep]
\item \textbf{Yüksek hassasiyet, küçük hacim (akademik etik birimi,
editor):} \emph{GPT-4o} (F1$=0.911$, $\$0.0020/$karar). Gerçek bir
referansi ``supheli'' yanlislamasi (false positive) yüksek olduundan, model
çıktıları insan denetimine \emph{aday liste} olarak kullanilmali.
\item \textbf{Yüksek hacim, bütçe kısıtlı (yayınevi pipeline'i):}
\emph{GPT-4o-mini} (F1$=0.862$, $\$0.0003/$karar). 8 kat daha ucuz, sadece
$0.05$ F1 puan kaybi. 100 bin referansli bir taramada $\$50$ yerine $\$6$.
\item \textbf{On-premise / hassas veri (kurumsal, gizli ar-ge):}
\emph{Llama 3.1-8B + RAG} (F1$=0.792$, $\$0$/karar). API'ye veri
sizmadan, GPT-4o-mini'nin \%92 performans bandında calisir; CrossRef
korpusunun guncellenmesi sorumluluğu kurumda olur.
\item \textbf{Hızlı prototip / 16~GB kısıtlı iş istasyonu:} Apple M2~Pro
yerel çalıştırma yeterlidir; p50 latency $\sim 2$~sn'dir.
\item \textbf{Kullanılması onerilmeyen:} \emph{RAG'siz Qwen2.5-7B}
(F1$=0.071$) -- modeli yalnızca RAG ile kullanmak gerekir, aksi takdirde
neredeyse rastgele karar verir.
\end{itemize}

\subsection{RAG katkısı: dengeli bir okuma}
\paragraph{Büyük ölçekli kazanım, beklenen yerlerde.} En büyük RAG kazanci
\emph{H1} (var olmayan referans) kategorisindedir: Llama için $+0.40$,
Qwen için pratik olarak $0 \to$ kısmı tespit. Bu, RAG'in birincil
mekanizmasiyla tam tutarlidir: bilgi tabaninda eslesme yoksa modele bu
sinyal aktarilir. \emph{H2} (metadata) kategorisi de ikinci büyük fayda
alanı olmuştur, cunku kutuphanedeki gerçek metadata refers'in bozulmuş
metadata'siyla karşı karsiya konulur.

\paragraph{Beklenmedik: tıp atfında gerileme.} Tablo~\ref{tab:catf1}'de
RAG'siz Llama'nın H4 F1$=0.483$ olmasi, RAG'li versiyonda $0.095$'e
düşmüştür (ayrıntılı rakamlar Bölüm~\ref{sec:sonuç} altında). Bu, i̇kili
\emph{tespit}in doğru, ancak \emph{tıp atfının} kayması anlamina gelir:
model ``evet bu bir halüsinasyon'' diyor, ama H4 (yıl hatası) yerine
\emph{H1} (uydurma) etiketi atiyor. Retrieval ipucu, modelin akıl yurutmesini
``eslesme arama'' moduna kanalize ediyor; ``yıl tutarsız'' veya ``bağlam
uyumsuz'' alt-sinyalleri zayifliyor. Bu, prompt mühendisliği yoluyla
hafifletilebilir: chain-of-thought \cite{wei2022cot} sablonunda modele
``once tipi listele, sonra karar ver'' yönergesi vermek, gelecek çalışma
için açık bir hat.

\paragraph{Pareto-onucu olarak ``açık kaynak + RAG''.} F1-maliyet
diyagramında RAG'li Llama (F1$=0.792$, \$0/karar) Pareto-onucu bolgede
yer alır: ondan daha ucuz hicbir model yoktur, ondan daha iyi yalnızca
ücretli kapalı kaynak modeller vardır. Bu, projenin \emph{somut praktik
katkısı}dir.

\subsection{Prompt mühendisliği gozlemi}
Pilot deneyde Qwen2.5-7B'in düşük recall'u, modelin referans
\emph{biçim}sel gecerliligini \emph{varlık} ile eşitlediğini göstermiştir
(parametric memory--only). Tam veri kümesinde bu kalibrasyon \emph{yo\u
gunlasarak} korunmuş, F1$=0.071$ ile aşağı yuvarlanmistir. RAG entegrasyonu
($+0.41$ F1) bu kalibrasyonu kısmen kirsa da, ek olarak chain-of-thought
\cite{wei2022cot} ve self-consistency \cite{wang2023selfconsistency}
teknikleri --- özellikle ``once tıp listele'' tipi yonergeler --- gelecek
çalışmada eklenebilir.

\subsection{Prompt mühendisliği gozlemi}
Pilot deneyde Qwen2.5-7B'in düşük recall'u, modelin referans
\emph{biçim}sel gecerliligini \emph{varlık} ile eşitlediğini göstermiştir
(parametric memory--only). Tam veri kümesinde bu kalibrasyonun korunması
durumunda, chain-of-thought \cite{wei2022cot} ve self-consistency
\cite{wang2023selfconsistency} teknikleri ek bir prompt-mühendisliği
ablasyonu olarak denenebilir.

\subsection{Sınırlamalar}
\begin{itemize}[nosep]
\item Benchmark İngilizce ve makine öğrenmesi/siber güvenlik alanında
yogunlasmistir; baska disiplinlerde genelleme tartışmalıdır.
\item RAG korpusu CrossRef ile siniridir; OpenAlex anahtarı abstract
zenginligine katkı yapar ancak gerekli degildir.
\item Modeller \texttt{temperature=0} ile deterministik ayardadir; sıcaklık
arttikca tutarlılık dusebilir (gelecek caliisma).
\item Beş kategori, tüm bibliyografya hatalarini kapsamaz (orn.\
sayfa-numarasi hatası, doi-yazım hatası); H2 altı sınıf olarak ele alinabilir.
\end{itemize}

% --- 7. Sonuç ve Gelecek Çalışmalar -----------------------------------------
\section{Sonuç ve Gelecek Çalışmalar}
\label{sec:gelecek}

\paragraph{Genel bulgu.}
Beş halüsinasyon türünü (H1--H5) ayırt eden 300 örnekli özgün benchmark'imiz
üzerinde dört LLM karşılaştırıldığında, kapalı kaynak GPT-4o en yüksek
performansi (F1$=0.911$) sergilemis; aynı sağlayıcının küçük modeli
GPT-4o-mini ise yaklaşık 1/8 maliyette F1$=0.862$ ile rekabetci kalmıştır.
Açık kaynak yerel modeller (Llama 3.1-8B, Qwen2.5-7B) parametric memory
tek başına yetersizdir (sirasiyla 0.560 ve 0.071); ancak CrossRef tabanlı
bir RAG katmanıyla güçlendirildiklerinde Llama F1$=0.792$ ile GPT-4o-mini'ye
yaklaşmış, Qwen ise $0.07 \to 0.48$ siçrayisi yapmistir ($\Delta$F1
ortalama $+0.32$, McNemar paired $p<0.001$). Bibliyografya halüsinasyonu
tespitinde \emph{retrieval, küçük açık kaynak modelleri pratik olarak
kullanilabilir hale getiren temel bilesendir}; veri gizliligi onemli
oldugunda RAG'li yerel modeller, kapalı kaynak modellerin yaklaşık \%87
performans bandında \emph{sıfır API maliyetiyle} çalışmaktadır.

\paragraph{Gelecek çalışmalar.}
\begin{itemize}[nosep]
\item \textbf{Fine-tuning (LoRA/QLoRA).} 500+ örnekli H1--H5 etiketli veriyle
açık kaynak modellerin ince ayarlanması (bonus, opsiyonel).
\item \textbf{Çok dilli benchmark.} Türkçe, Almanca, Çince akademik atıflara
genisleme.
\item \textbf{Aktif retrieval.} Modelin hangi alanlarda RAG çağrısı yapmasi
gerektiğine kendisinin karar verdigi adaptif retrieval.
\item \textbf{Insan-makine isbirligi.} Sistemin ``supheli'' işaretlediği
örnekleri sırayla insan denetimine getiren akıllı kuyruk.
\end{itemize}

% --- Kaynakça ---------------------------------------------------------------
\printbibliography[heading=bibnumbered, title={Kaynakça}]

% =============================================================================
\appendix

\section{YZ Araçları Kullanım Notu}
\label{appx:ai}
Bu proje sırasında kullanılan YZ araçları, aşamaları ve sorumluluk paylaşımı:

\begin{tabularx}{\columnwidth}{l l X}
\toprule
\textbf{Araç} & \textbf{Aşama} & \textbf{Amaç} \\
\midrule
Claude (Anthropic) & Kod, rapor, sunum & Repository iskeleti, rapor metni düzenleme, sunum slaytlarının hazırlanması \\
ChatGPT (OpenAI)   & Metin             & Bazı paragraf revizyonları ve dilbilgisi kontrolü \\
GPT-4o, GPT-4o-mini (OpenAI API) & Deney & Benchmark üzerinde değerlendirilen closed-source modeller \\
Llama 3.1-8B, Qwen2.5-7B (Ollama) & Deney & Benchmark üzerinde değerlendirilen open-source modeller \\
\bottomrule
\end{tabularx}

\noindent Benchmark anotasyonu, deney çalıştırma, hata analizi, raporun nihai metni
ve sunumun nihai içeriği sorumluluğu öğrenciye aittir. YZ otomatik
etiketlemesi ground-truth olarak kullanılmamıştır.

\section{Prompt Şablonu}
\label{appx:prompt}
\textbf{Sistem mesaji:} ``Sen akademik referans doğrulama uzmanisin. Verilen
referansın gerçek bir akademik yayın olup olmadığını ve metadata tutarliligini
analiz et. Su halüsinasyon türlerini ayırt et: H1 = tamamen var olmayan
referans; H2 = metadata halüsinasyonu; H3 = semantik uyumsuzluk;
H4 = kronolojik halüsinasyon; none = doğru referans. Yalnızca su JSON
semasiyla yanıt ver:
\texttt{\{"hallucination": true/false, "confidence": 0.0-1.0,
"type": "H1|H2|H3|H4|none", "reason": "kisa gerekçe"\}}.''

\noindent\textbf{Kullanıcı mesaji:} \texttt{REFERANS:\\<girdi>\\\\(varsa)\
ATIF BAĞLAMI:\\<bağlam>\\\\(varsa)\ BILGI TABANI (RAG):\\<getirilen kayıtlar>\\\
Bu referansi analiz et ve JSON karari ver.}

\section{Benchmark Örnek Girdiler}
\label{appx:benchmark}
Her kategoriden bir örnek (\texttt{data/benchmark/benchmark.json}):

\noindent\textbf{ID 1 -- H1 -- orta zorluk.}\\
\texttt{D. Anderson, J. Halloran, H. Esposito, D. Vasquez, ``Privacy-Preserving Zero-Day Exploit Forecasting,'' in Proc.\ USENIX Security Symposium, 2024.}\\
\textit{GT}: hallucination$=$true, type$=$H1.
\textit{Gerekçe}: Tamamen uydurma referans; bu başlık/yazar kombinasyonuna
sahip yayimlanmis bir makale yoktur.

\medskip
\noindent\textbf{ID 81 -- H2 -- orta zorluk.}\\
\texttt{T.~A.~Tang, L.~Mhamdi, D.~McLernon et al., ``Deep learning approach for Network Intrusion Detection in Software Defined Networking,'' in [yer], 2019. doi: 10.1109/...}\\
\textit{GT}: hallucination$=$true, type$=$H2.
\textit{Gerekçe}: Yayın yılı gerçek kayittan farklı; metadata uyumsuzluğu.

\medskip
\noindent\textbf{ID 141 -- H3 -- kolay zorluk.}\\
\texttt{H.~H.~Aghdam, E.~J.~Heravi, D.~Puig, ``Explaining Adversarial Examples by Local Properties of Convolutional Neural Networks,'' in Proc.\ Int.\ Conf., 2017.}\\
\textit{Bağlam}: \emph{``This study evaluates the carbon-sequestration capacity of mangrove forests across three coastal regions.''}\\
\textit{GT}: hallucination$=$true, type$=$H3.
\textit{Gerekçe}: Referans gerçek olsa da atıf bağlamı (mangrov ormanlari)
ile semantik uyumsuzluk.

\medskip
\noindent\textbf{ID 201 -- H4 -- orta zorluk.}\\
\texttt{P.-Y.~Chen, H.~Zhang, Y.~Sharma et al., ``ZOO: Zeroth Order Optimization Based Black-Box Attacks,'' in [yer], 1995.}\\
\textit{GT}: hallucination$=$true, type$=$H4.
\textit{Gerekçe}: 1995 tarihi imkânsız; ZOO 2017'de yayimlanmistir
(kronolojik çelişki).

\medskip
\noindent\textbf{ID 241 -- H5 -- orta zorluk.}\\
\texttt{B.~M.~Khammas, ``Ransomware Detection using Random Forest Technique,'' in ICT Express, 2020. doi: 10.1016/j.icte.2020.11.001.}\\
\textit{GT}: hallucination$=$false (negatif).
\textit{Gerekçe}: Tüm alanlar (başlık, yazar, yıl, yer, DOI) CrossRef'te
doğrulanmıştır.

\section{Tekrarlanabilirlik Notu}
\label{appx:repro}
Tüm kod \url{https://github.com/novruzamirov/bgk601-bibhallu} adresinde MIT
lisansıyla yayınlanmıştır. Smoke-test (internet/GPU/anahtar gerekmez):
\texttt{bash scripts/smoke\_test.sh}. Tam pipeline için:
\texttt{scripts/0\_setup.sh} $\to$ \texttt{1\_build\_dataset.sh} $\to$
\texttt{2\_run\_eval.sh} $\to$ \texttt{3\_build\_rag.sh} $\to$
\texttt{4\_rag\_ablation.sh} $\to$ \texttt{5\_analyze.sh}.

\end{document}
